% ===== PRÉAMBULE CANONIQUE — à coller VERBATIM en tête de CHAQUE .tex =====
\documentclass[11pt,a4paper]{article}
\usepackage[utf8]{inputenc}\usepackage[T1]{fontenc}\usepackage{lmodern}
\usepackage[french]{babel}
\usepackage{amsmath,amssymb,mathtools}\usepackage{icomma}
\usepackage[version=4]{mhchem}          % \ce{...} : équations & formules
\usepackage{siunitx}\sisetup{locale=FR} % \qty{}{}, \si{}, \ang{}
\usepackage{chemfig}                    % \chemfig{...} : structures développées
\usepackage{xcolor}\usepackage{colortbl}
\usepackage{tikz}\usepackage{pgfplots}\pgfplotsset{compat=1.18}
\usepackage[most]{tcolorbox}\usepackage{enumitem}\usepackage{array,booktabs}
\usepackage[a4paper,margin=2.2cm]{geometry}
\usetikzlibrary{arrows.meta,calc,angles,quotes,positioning,patterns,backgrounds,shapes.geometric}
\definecolor{primary}{RGB}{222,49,99}\definecolor{primarydark}{RGB}{150,22,55}
\definecolor{defcol}{RGB}{0,114,178}\definecolor{methcol}{RGB}{52,92,148}
\definecolor{excol}{RGB}{0,135,90}\definecolor{attncol}{RGB}{200,40,55}
\tcbset{boxbase/.style={enhanced,breakable,boxrule=0.9pt,arc=2.5pt,
  left=3mm,right=3mm,top=2.3mm,bottom=2.3mm,fonttitle=\bfseries,coltitle=white}}
% --- boîtes TUTORIEL (noms EXACTS, ne pas renommer) ---
\newtcolorbox{cle}[1][]{boxbase,colback=defcol!6,colframe=defcol,
  title={\textsc{À retenir}\ifx\relax#1\relax\relax\else\textmd{ : \itshape #1}\fi}}
\newtcolorbox{methode}[1][]{boxbase,colback=methcol!7,colframe=methcol,sharp corners,
  title={\textsc{Méthode}\ifx\relax#1\relax\relax\else\textmd{ --- \itshape #1}\fi}}
\newtcolorbox{exemplebox}[1][]{boxbase,colback=excol!5,colframe=excol,
  title={\textsc{Exemple}\ifx\relax#1\relax\relax\else\textmd{ : \itshape #1}\fi}}
\newtcolorbox{piege}[1][]{boxbase,colback=attncol!6,colframe=attncol,
  title={\textsc{Piège}\ifx\relax#1\relax\relax\else\textmd{ : \itshape #1}\fi}}
\newtcolorbox{remarque}[1][]{boxbase,colback=black!4,colframe=black!45,
  title={\textsc{Remarque}\ifx\relax#1\relax\relax\else\textmd{ : \itshape #1}\fi}}
% --- boîtes EXERCICE / CORRIGÉ (noms EXACTS, auto-comptées) ---
\newtcolorbox[auto counter]{exo}[1][]{boxbase,colback=primary!4,colframe=primary,
  title={\textsc{Exercice}~\thetcbcounter\ifx\relax#1\relax\relax\else\textmd{ --- \itshape #1}\fi}}
\newtcolorbox[auto counter]{corrige}[1][]{boxbase,colback=excol!4,colframe=excol,
  title={\textsc{Corrigé}~\thetcbcounter\ifx\relax#1\relax\relax\else\textmd{ --- \itshape #1}\fi}}
\newcommand{\R}{\mathbb{R}}\newcommand{\N}{\mathbb{N}}\newcommand{\Z}{\mathbb{Z}}\newcommand{\ds}{\displaystyle}
% (un \usepackage{fancyhdr}+en-tête est possible ; il est ignoré côté web)
% =========================================================================
\begin{document}
% THEME_TITLE: Les oxydes (métalliques et non métalliques)

\begin{cle}[un oxyde, deux règles selon l'autre élément]
Un \textbf{oxyde} est un composé binaire contenant de l'\textbf{oxygène} (au n.o.\ $-\mathrm{II}$). La règle
de nomenclature dépend de la nature de l'\emph{autre} élément :
\begin{itemize}[leftmargin=1.5em,topsep=2pt,itemsep=1pt]
  \item \textbf{métal} $\to$ \emph{oxyde métallique} : on nomme avec la \textbf{valence} du métal ;
  \item \textbf{non-métal} $\to$ \emph{oxyde non métallique} : on nomme avec un \textbf{préfixe multiplicateur}.
\end{itemize}
\end{cle}

\begin{methode}[{oxyde métallique : « oxyde de [métal] (valence) »}]
On écrit \textbf{« oxyde de [nom du métal] »}, en précisant la valence en chiffres romains
\emph{uniquement} si le métal peut en prendre plusieurs (\ce{Fe}, \ce{Cu}, \ce{Sn}, \ce{Pb}, \ce{Hg}).
Pour un métal à valence \emph{unique} (\ce{Na}, \ce{Ca}, \ce{Mg}, \ce{Al}, \ce{Zn}, \ce{Ag}, \ce{K}), on
ne l'indique pas. Pour écrire la formule, on croise la valence du métal avec $2$ (celle de \ce{O^2-}).
\end{methode}

\begin{methode}[oxyde non métallique : préfixe $=$ rapport O/X]
On place devant « oxyde » un préfixe traduisant le rapport $\dfrac{\text{nombre d'oxygènes}}{\text{nombre d'atomes }X}$,
puis « de [nom de $X$] ».
\begin{center}\footnotesize
\renewcommand{\arraystretch}{1.05}
\begin{tabular}{cl@{\qquad}cl@{\qquad}cl}
\toprule
O/X & préfixe & O/X & préfixe & O/X & préfixe \\
\midrule
1/2 & hémi   & 3/1 & tri   & 3/2 & sesqui \\
1/1 & mono   & 4/1 & tétra & 5/2 & hémipent \\
2/1 & di     & 5/1 & pent  & 7/2 & hémihept \\
\bottomrule
\end{tabular}
\end{center}
Quand $X$ est présent \emph{une} fois (\ce{XO_n}), le préfixe est simplement le nombre d'oxygènes
(\ce{CO}=mono, \ce{NO2}=di, \ce{SO3}=tri). Quand $X$ est présent \emph{deux} fois (\ce{X2O_n}), le rapport
vaut $n/2$ et l'on emploie les préfixes fractionnaires (\ce{N2O3}=sesqui, \ce{N2O5}=hémipent).
\end{methode}

\begin{exemplebox}[lire sept oxydes]
\ce{CaO} : oxyde de calcium ; \ce{FeO} : oxyde de fer(II) ; \ce{Fe2O3} : oxyde de fer(III) ;
\ce{Ag2O} : oxyde d'argent ; \ce{CO} : monoxyde de carbone ; \ce{NO2} : dioxyde d'azote ;
\ce{Cl2O5} : hémipentoxyde de chlore.
\end{exemplebox}

\begin{piege}[les erreurs qui se répètent au concours]
\begin{itemize}[leftmargin=1.4em,topsep=2pt,itemsep=1pt]
  \item \ce{CuO} $=$ oxyde de cuivre\textbf{(II)} et \ce{Cu2O} $=$ oxyde de cuivre\textbf{(I)} :
        ne pas intervertir les valences.
  \item pour un métal à valence unique, \emph{ne pas} écrire la valence (\ce{ZnO} $=$ oxyde de zinc, pas « (II) »).
  \item pour un oxyde non métallique, c'est le \textbf{rapport} O/X qui compte, pas le seul nombre d'oxygènes :
        \ce{N2O5} est un \emph{hémipent}oxyde ($5/2$), pas un « pentoxyde ».
\end{itemize}
\end{piege}

\begin{remarque}[une autre nomenclature des oxydes non métalliques]
On rencontre aussi la nomenclature \emph{de stœchiométrie} qui compte les \emph{deux} éléments :
\ce{N2O} y devient « mon\emph{oxyde} de di\emph{azote} » et \ce{NO2} « dioxyde d'azote ». Dans cette fiche,
on suit la convention du \textbf{rapport O/X} : \ce{N2O} $=$ hémioxyde d'azote ($1/2$). Sache reconnaître les
deux écritures.
\end{remarque}

\end{document}
